\cotefile{33}
\scanpage{18}
à des $f$ de type fini au lieu de '$f$' propre. Se ramène au cas d'une \uline{imm. ouverte}. Se résoud par l'\unsure{affirmative|affirm.} lorsqu'on dispose de la résolution des singularités\inserted{$+$ pureté}, p. ex. au cas où \inserted{pour ``bons préschémas''}, \unsure{ou plus|ou plan|ou plutôt}\inserted{\unsure{bons présch|bien précisé}} \unsure{cf. le|cgt le|gén.} cas lorsqu'on disposera de résolution des singularités \deleted{\unsure{des|de}} dans ce cadre....

14) \uline{Les} $R^{i}_{!}f\unsure{{}_{*}|{}}(F)$.

Soit $f:X\longrightarrow Y$ morphisme compactifiable, i.e. qui se factorise en

\diagram{\begin{tikzcd}[ampersand replacement=\&,column sep=large]
X \arrow[r,dash,"i","{\mathrm{imm.\ ouv.}}"'] \& \bar X \arrow[r,"{\bar f}","{\mathrm{propre}}"'] \& Y
\end{tikzcd}}

(p. ex. $f$ qu. proj. et $\deleted{f}Y$ quasi-cpt\unsure{?|.}).

Pour tt faisceau de torsion $F$ sur $X$, considérons $i_!(F)$ [exact en $F$] puis les $R^{q}_{\deleted{\unsure{q|i}}}\bar f_*(i_!(F))$, forment un foncteur coh. en $F$. Il résulte de façon quasi-triviale du th. de changement de base que ce foncteur coh.

% archive sheet number: 17
\scanpage{19}
\mspage{- 191 -}\hfill n\textsuperscript{o} 391

racines de $R$ sont telles que $(u\mid u)=1$ ou $(u\mid u)=2$, (en excluant le cas des systèmes de racines de type $G_2$, dont la structure a été déterminée au \S~6, n\textsuperscript{o}~12).

Inversement, partons d'un réseau $L$ dans un espace euclidien $V$, et soit $R$ la partie formée de \uline{tous} les vecteurs de $L$ tels que $(u\mid u)=\lambda$, (ou vérifiant l'une des deux conditions $(u\mid u)=\lambda$, $(u\mid u)=2\lambda$), pour un $\lambda$ convenable. Alors, si $R$ engendre $V$ et si on suppose de plus que pour tout couple $(u,v)$ de vecteurs de $R$, le nombre $2(u\mid v)/(u\mid u)$ est \uline{entier}, $R$ est un \uline{système de racines} : en effet, si $s_u$ est la réflexion orthogonale de vecteur $u$, $s_u(v)$ appartient alors au sous-groupe de $V$ engendré par $u$ et $v$, donc \uline{a fortiori} à $L$, et en outre $s_u(v)$ a même norme que $v$, donc est dans $R$, ce qui prouve $(\mathrm{SR}_{\mathrm I})$ ; comme $(\mathrm{SR}_{\mathrm{II}})$ et $(\mathrm{SR}_{\mathrm{III}})$ sont \unsure{vérifiés|vérifiées} par hypothèse, cela prouve notre assertion. Il restera naturellement à vérifier dans chaque cas que le système $R$ est réduit et irréductible.

Nous allons voir qu'il est possible d'obtenir tous les systèmes de racines correspondant aux schémas de Dynkin du th. 2, ($G_2$ exclus), en prenant pour $V$ un espace numérique $R^n$, dont nous désignerons par $\varepsilon_i$, $(1\leq i\leq n)$, la base canonique, ($(\varepsilon_i^*)$ étant la base duale), et où la forme métrique admet $(\varepsilon_i)$ pour \uline{base orthonormale} ; on considèrera alors dans $V$ les réseaux $L_i$, $(i=0,1,2,3)$, suivants notés aussi $L_i(V)$, (et leurs intersections par certains sous-espaces de $V$) :

1) Le réseau $\unsure{L_o|L_0}$ est le $Z$-module de base $(\varepsilon_i)$ ; il est clair que $(u\mid v)$ est entier pour $u$, $v$ dans $L_o$. En outre, les vecteurs de $L_o$ tels que $(u\mid u)=1$ sont les vecteurs $\pm\varepsilon_i$, $(1\leq i\leq n)$ ; ceux pour lesquels $(u\mid u)=2$ sont les vecteurs $\pm\varepsilon_i\pm\varepsilon_j$ pour $i\neq j$.

% archive sheet number: 18
\scanpage{20}
à isom près est \uline{indépendant} du choix de la \unsure{fact.|fact\textsuperscript{n}} Notation $R^q_! f_*(F)$.

\[
R^0_! f_*(F)=f_!(F)=\text{image directe à support propre sur }Y.
\]

\deleted{Propriétés} Propriétés plus surprenantes à certains égards que $R^q f_*(F)$, car \uline{aussi bon que si $f$ propre}

(i) \uline{Comparaison avec cas classique} sur $\mathbb C$

(ii) \uline{Th. de changement de base} \deleted{\unsure{cf|cf.|c}}, \uline{dim} \uline{coh.}

(iii) \uline{\deleted{Suite spectrale de Leray}}\inserted{Th. de finitude}

(iv) Suite spectrale de Leray

% slanted
\margin{phénomènes spéciaux}
(v) Si $f$ quasi-fini\inserted{\unsure{sép.}}, alors les $R^i_! f_*(F)=0$

\deleted{D'autres \unsure{faits|points|part} =}
(vi) Si $f$ \uline{étale}, on a une trace can.

\[
\mathrm{Tr}_f(F):f_!(F)\unsure{\mathrel{-}|\longrightarrow}F
\]

[d'où loi covariante, pour $f:X'\mathrel{-}X$ morphisme étale -- p. ex. imm ouverte -- ] de préschémas compactifiables sur $Y$

\[
R^q_!g'_*(F')\mathrel{-}R^q_!g_*(f_!(F'))
\]

ou encore

\[
R^q_!g'_*(\deleted{\unsure{F}}f^*(F))\longrightarrow R^q_!g_*(F)\qquad\big]
\]

(vi) \uline{Suite exacte de cohomologie relative}

% archive sheet number: 19
\scanpage{21}
\mspage{- 189 -}\hfill n\textsuperscript{o} 391

En effet, (\S~6, n\textsuperscript{o}~5, prop. 12), ce groupe quotient est isomorphe au sous-groupe $\unsure{G_o|G_0}$ de $A(R)$ laissant stable un système de racines simples $S(C)$ de $R$ ; il est clair que tout élément de $G_o$ détermine un automorphisme unique du schéma de Dynkin de $R$ relatif à $C$ ; mais inversement, un tel automorphisme se prolonge en un automorphisme unique de $R$ en vertu de la prop. 2, d'où le corollaire.

Les remarques du début et les définitions précédentes montrent que le schéma de Dynkin d'un système de racines $R$, relatif à un système de racines simples $(v_i)_{1\leq i\leq\ell}$, détermine le schéma de Coxeter du groupe de Weyl $W(R)$ de la manière suivante : les graphes sous-jacents sont les mêmes ; en outre, si $\{v_i,v_j\}$ est une arête du schéma de Dynkin, le nombre $m_{ij}$ est égal à 3 si $f(v_i)=f(v_j)$, à 4 si $f(v_i)/f(v_j)$ est égal à $\frac{1}{2}$ ou à 2, à 6 si $f(v_i)/f(v_j)$ est égal à $1/3$ ou à 3. On verra par contre ci-dessous qu'un même schéma de Coxeter peut correspondre à deux schémas de Dynkin non isomorphes, (en fait, il ne peut correspondre qu'à \uline{deux} tels schémas au plus, et cela ne se produit que pour les schémas de Coxeter de type $B_\ell$ ; pour tous les autres types, le schéma de Coxeter de $W(R)$ détermine le schéma de Dynkin de $R$ à isomorphisme près).

THEOREME 2.- \uline{Si} $R$ \uline{est un système de racines réduit irréductible}, \uline{son schéma de Dynkin est nécessairement isomorphe à l'un des suivants}, (\uline{qui sont deux à deux non isomorphes}) :

\diagram{\begin{tikzcd}[ampersand replacement=\&,column sep=1.1em,row sep=.35em,cells={nodes={inner sep=0pt}},every arrow/.append style={start anchor=center,end anchor=center}]
A_\ell \& \vert \arrow[r,dash] \& \vert \arrow[r,dash] \& \vert \arrow[r,dash,dotted] \& \vert \arrow[r,dash] \& \vert \& (\ell\geq 1~\text{\uline{sommets}})
\end{tikzcd}}

\diagram{\begin{tikzcd}[ampersand replacement=\&,column sep=1.1em,row sep=.35em,cells={nodes={inner sep=0pt}},every arrow/.append style={start anchor=center,end anchor=center}]
B_\ell \& |[label=above:{$2$}]|\vert \arrow[r,dash] \& |[label=above:{$2$}]|\vert \arrow[r,dash] \& |[label=above:{$2$}]|\vert \arrow[r,dash,dotted] \& |[label=above:{$2$}]|\vert \arrow[r,dash] \& |[label=above:{$2$}]|\vert \arrow[r,dash] \& |[label=above:{$1$}]|\vert \& (\ell\geq 2~\text{\uline{sommets}})
\end{tikzcd}}

\diagram{\begin{tikzcd}[ampersand replacement=\&,column sep=1.1em,row sep=.35em,cells={nodes={inner sep=0pt}},every arrow/.append style={start anchor=center,end anchor=center}]
C_\ell \& |[label=above:{$1$}]|\vert \arrow[r,dash] \& |[label=above:{$1$}]|\vert \arrow[r,dash] \& |[label=above:{$1$}]|\vert \arrow[r,dash,dotted] \& |[label=above:{$1$}]|\vert \arrow[r,dash] \& |[label=above:{$1$}]|\vert \arrow[r,dash] \& |[label=above:{$2$}]|\vert \& (\ell\geq 3~\text{\uline{sommets}})
\end{tikzcd}}

\diagram{\begin{tikzcd}[ampersand replacement=\&,column sep=1.1em,row sep=.35em,cells={nodes={inner sep=0pt}},every arrow/.append style={start anchor=center,end anchor=center}]
 \&  \&  \&  \&  \&  \&  \& \diagdown \\
D_\ell \& \vert \arrow[r,dash] \& \vert \arrow[r,dash] \& \vert \arrow[r,dash,dotted] \& \vert \arrow[r,dash] \& \vert \arrow[r,dash] \& {} \arrow[ur,dash]\arrow[dr,dash] \&  \& (\ell\geq 4~\text{\uline{sommets}}) \\
 \&  \&  \&  \&  \&  \&  \& \diagup
\end{tikzcd}}

\diagram{\begin{tikzcd}[ampersand replacement=\&,column sep=1.1em,row sep=.35em,cells={nodes={inner sep=0pt}},every arrow/.append style={start anchor=center,end anchor=center}]
E_6 \& \vert \arrow[r,dash] \& \vert \arrow[r,dash] \& \vert \arrow[r,dash] \arrow[d,dash] \& \vert \arrow[r,dash] \& \vert \\
 \&  \&  \& -
\end{tikzcd}}

% archive sheet number: 20
\scanpage{22}
\mspage{\textcircled{15}}
On fixe entier $n$ premier aux car. résiduelles

\uline{Classe fondamentale locale}

Supposons $Y\hookrightarrow X$ ``intersection complète de dim relative $1$'' et supposons d'abord $Y=V(\varphi)$, $\varphi$ section régulière de $\mathcal O_X$. Alors si $\unsure{u|U}:X-Y=U\longrightarrow X$, donc \deleted{$i_*(\underline O_U)\sim O_X[\frac{1}{\varphi}]$}, \deleted{$i_*(\underline O^*_U)$ et} $\varphi\in\Gamma(i_*(\deleted{\unsure{\underline O}}\underline G_{mU}))$ \deleted{et $\unsure{\sim|=|\to}\ \in\Gamma(\unsure{X;|i_*}\,i_*(\unsure{G|O}^*_{mU}))$} \deleted{D'\unsure{où|ai|au}} ailleurs $\varphi$ est \unsure{déterminé|déterminée} \inserted{par $Y$} mod multiplication par un $\in\Gamma(X,i_*(\deleted{\unsure{\underline O}}G_{mX}))$, donc \deleted{l'image} $\varphi$ définit une section \deleted{$\varphi\in$} \uline{canonique} de \deleted{$i_*(\underline O^*_U)/$} $i_*(G_{mU})/G_{mX}$. Or on a la suite exacte canonique

\diagram{\begin{tikzcd}[ampersand replacement=\&,column sep=.75em,row sep=.5em]
0 \arrow[r] \& \underline H^0_Y(G_{mX}) \arrow[r]\arrow[d,equal] \& G_{mX} \arrow[r,dash] \& i_*(G_{mU}) \arrow[r,dash,"d"] \& \underline H^1_Y(G_{mX}) \arrow[d] \\
\& 0 \& \& \& 0
\end{tikzcd}}

qui donne ici

\[
\deleted{\unsure{H^1|H}}\ \underline H^1_Y(G_{mX})\xrightarrow{\sim}i_*(G_{mU})/\mathrm{Im}\,G_{mX}\quad \unsure{:|;})
\]

et montre que $d\varphi$ est une section \uline{canonique} de $\underline H^1_Y(G_{mX})$, que nous allons noter $\deleted{\unsure{\xi|\gamma|\varepsilon}_{Y/X}}\underline\beta_{Y/X}$. D'autre part, on a \deleted{une}\inserted{la} suite exacte de Kummer

\[
\unsure{0|\cdot}\mathrel{-}\unsure{\mu_{n_X}|\mu_{nX}}\longrightarrow G_{mX}\overset{n}{\mathrel{-}} G_{mX}\longrightarrow 0
\]

% archive sheet number: 21
\scanpage{23}
\mspage{- 218 -}\hfill n\textsuperscript{o} 391

(VII) La demi-somme des racines positives est la somme des poids fondamentaux, (\S~6, n\textsuperscript{o}~9, prop. 8), et vaut donc :

\[
\begin{aligned}
\rho&=\varepsilon_2+2\varepsilon_3+3\varepsilon_4+4\varepsilon_5+4(\varepsilon_8-\varepsilon_7-\varepsilon_6)=\\
&=8\alpha_1+11\alpha_2+15\alpha_3+21\alpha_4+15\alpha_5+8\alpha_6.
\end{aligned}
\]

(VIII) De même qu'au n\textsuperscript{o}~11, le réseau des vocables $\unsure{\Gamma_o|\Gamma_0}(R'')$ est l'associé \uline{dans} $V''$ du réseau $\Gamma_2^*$ de $V''{}^*$ \uline{projection orthogonale} dans $V''{}^*$ de $\Gamma_o(R^*)$, (lorsque $V''{}^*$ est identifié au sous-espace de $V^*$ orthogonal à $\varepsilon_7^*+\varepsilon_8^*$ et $\varepsilon_7^*-\varepsilon_6^*$). Le groupe $\Gamma(R'')/\Gamma_o(R'')$ est donc \uline{isomorphe au groupe} $\Gamma_2^*/\Gamma_o(R''{}^*)$, qu'il reste à calculer. Or, les projections orthogonales de $\varepsilon_8^*$, $-\varepsilon_7^*$ et $-\varepsilon_6^*$ sur $V''{}^*$ sont toutes trois égales $\frac{1}{3}(\varepsilon_8^*-\varepsilon_7^*-\varepsilon_6^*)$ et les autres $\varepsilon_i^*$ sont dans $V''{}^*$ ; on en déduit aisément que $\Gamma_2^*$ est engendré par $\Gamma_o(R''{}^*)$ et par l'unique vecteur $\mu^*=\frac{1}{3}(\varepsilon_8^*-\varepsilon_7^*-\varepsilon_6^*)$ ; mais comme $3\mu^*\in\Gamma_o(R''{}^*)$, on voit que $\Gamma_2^*/\Gamma_o(R''{}^*)$ est isomorphe à $Z/3Z$ et par suite $\Gamma(R'')/\Gamma_o(R'')$ est \uline{isomorphe à} $Z/3Z$, et l'indice de connexion $f=3$.

(IX) et (X). Le calcul de la plus grande racine et de l'indice de connexion permet, grâce au \S~6, n\textsuperscript{o}~10, prop. 23, de calculer l'ordre du groupe de Weyl $W(R'')$, qui est égal à :

\[
6!1.2.2.3.2.1.3=2^7.3^4.5=51.840\unsure{.|:}
\]

La suite des exposants a 6 termes compris entre 1 et 11, et comprend les termes 1, 5, 7, 11 qui sont étrangers à 12, (\S~8, n\textsuperscript{o}~5, prop. 6). Les deux derniers termes de la suite $m$, $m'$ sont des entiers tels que :

\[
\begin{aligned}
m+m'&=12\\
(m+1)(m'+1)(1+1)(5+1)(7+1)(12+1)&=2^7.3^4.5
\end{aligned}
\]

% archive sheet number: 22
\scanpage{24}
d'où

\[
\partial_n:\underline H^1_Y(G_{mX})\longrightarrow\underline H^2_Y(\unsure{\mu_{n_X}|\mu_{nX}})
\]

et nous poserons

\[
\deleted{\unsure{\xi|\varphi}}\ \underline\gamma^{(n)}_{Y/X}=\partial_n(\underline\beta_{Y/X})\quad\in\Gamma(X,\underline H^2_Y(\mu_{nX}))
\]

en laissant tomber \unsure{le|les} exposant $n$ s'il n'y a pas d'ambiguité possible.

\inserted{La définition de} Cette section \deleted{satisfait} commute à la localisation de Zariski (et $=$ à tt \unsure{changt|changemt} de base $X'\longrightarrow X$ tel que $Y'$ reste une ``int. complète'' i.e. $\varphi'$ reste section régulière) ce qui permet de définir $\gamma^{(n)}_{Y/X}$ en général (sans équation globale). Si \inserted{maintenant} $Y$ est intersection complète de codim $d$, alors localemt $Y$ est de la forme

\[
Y=Y_1\cap\ldots\cap Y_d
\]

et on posera

\[
\underline\gamma_{Y/X}=\underline\gamma^{\deleted{\unsure{(n)}}}_{Y_1/X}\mathbin{.}\ \mathbin{.}\underline\gamma_{Y_d/X}\quad\in\Gamma(Y,\underline H^{2d}_Y(\mu_n^{\otimes d}))
\]

(\uline{cup-produit}), et on prouve que cette section ne dépend pas du choix des $Y_i$, donc se définit globalement.

Propriétés générales :

% archive sheet number: 23
\scanpage{25}
\mspage{- 217 -}\hfill n\textsuperscript{o} 391

(IV) La racine $\frac{1}{2}(\varepsilon_1+\varepsilon_2+\varepsilon_3+\varepsilon_4+\varepsilon_5-\varepsilon_6-\varepsilon_7+\varepsilon_8)$ est la plus grande des racines de $R''$ pour l'ordre induit par celui considéré sur $R^8$, (n\textsuperscript{o}~10) ; c'est donc la plus grande racine $\tilde\alpha$, (\S~6, n\textsuperscript{o}~6, cor. 2 de la prop. 13) ; on a $\tilde\alpha=\alpha_1+2\alpha_2+2\alpha_3+3\alpha_4+2\alpha_5+\alpha_6$, et on vérifie que son ordre est $h-1=11$. Cette racine est orthogonale à toutes les racines simples $\alpha_i$, $(1\leq i\leq6)$, sauf $\alpha_2$, et l'on a $(\tilde\alpha\mid\alpha_2)=1$.

(V) Le même raisonnement qu'au n\textsuperscript{o}~11 montre que les \uline{coracines} de $R''$ s'obtiennent en remplaçant $\varepsilon_i$ par $\varepsilon_i^*$ dans les expressions des racines de $R''$, et $R''{}^*$ est isomorphe à $R''$. Toutes les racines ayant la même longueur, la \uline{forme bilinéaire canonique} pour $R$ est $(x\mid y)/2h=(x\mid y)/24$ et la longueur des racines est $1/2\sqrt{3}$, (\S~8, n\textsuperscript{o}~5, \uline{Remarque}). Le nombre des racines non orthogonales à une racine donnée est $4h-6=42$ et le nombre $\gamma$ est égal à $h^2=144$.

(VI) Le calcul des poids fondamentaux, formant la base duale de $(\alpha_i^*)$, $(1\leq i\leq6)$, dans $V''$, donne :

\[
\begin{aligned}
\omega_1&=\frac{2}{3}(\varepsilon_8-\varepsilon_7-\varepsilon_6)=\frac{1}{3}(4\alpha_1+3\alpha_2+5\alpha_3+6\alpha_4+4\alpha_5+2\alpha_6)\\
\omega_2&=\frac{1}{2}(\varepsilon_1+\varepsilon_2+\varepsilon_3+\varepsilon_4+\varepsilon_5-\varepsilon_6-\varepsilon_7+\varepsilon_8)=\\
&=\alpha_1+2\alpha_2+2\alpha_3+3\alpha_4+2\alpha_5+\alpha_6=\tilde\alpha\\
\omega_3&=\frac{5}{6}(\varepsilon_8-\varepsilon_7-\varepsilon_6)+\frac{1}{2}(-\varepsilon_1+\varepsilon_2+\varepsilon_3+\varepsilon_4+\varepsilon_5)=\\
&=\frac{1}{3}(5\alpha_1+6\alpha_2+10\alpha_3+12\alpha_4+8\alpha_5+4\alpha_6)\\
\omega_4&=\varepsilon_3+\varepsilon_4+\varepsilon_5-\varepsilon_6-\varepsilon_7+\varepsilon_8=2\alpha_1+3\alpha_2+4\alpha_3+6\alpha_4+4\alpha_5+2\alpha_6\\
\omega_5&=\frac{2}{3}(\varepsilon_8-\varepsilon_7-\varepsilon_6)+\varepsilon_4+\varepsilon_5=\\
&=\frac{1}{3}(4\alpha_1+6\alpha_2+8\alpha_3+12\alpha_4+10\alpha_5+5\alpha_6)\\
\omega_6&=\frac{1}{3}(\varepsilon_8-\varepsilon_7-\varepsilon_6)+\varepsilon_5=\\
&=\frac{1}{3}(2\alpha_1+3\alpha_2+4\alpha_3+6\alpha_4+5\alpha_5+4\alpha_6).
\end{aligned}
\]

% archive sheet number: 24
\scanpage{26}
% slanted
\margin{\uline{propriétés caractéristiques}}
\[
\left\{\begin{aligned}
&\text{(i) Compatibilité avec \unsure{chang|chg}\textsuperscript{t} de base admissible}\\
&\text{(ii) Si }\ Y=V(\varphi),\qquad \underline\gamma^{(n)}_{Y/X}=d(\beta^{(n)}_{Y/X})\ \text{comme dessus}\\
&\text{(iii) }\qquad \underline\gamma^{(n)}_{Y.Y'/X}=\underline\gamma^{(n)}_{Y/X}\mathbin{.}\underline\gamma^{(n)}_{Y'/X}
\end{aligned}\right.
\]
si $Y.Y'$ ``sans dim. excédentaire''

\[
\Downarrow
\]

(iv) Compatibilité pour $n$ variable

(\uline{NB} La ``bonne'' classe se trouve dans $\Gamma(Y,\underline H^{2d}_Y(\mathbb Z_\ell(d)))$, il y en a une pour tt $\ell$ premier aux car. résiduelles de $Y$).

\uline{Application à la construction de classes de coh. associées à certains cycles.}

Supposons $\underline H^i_Y((\mathbb Z/n\mathbb Z)_X)=0$ \unsure{pr|gr} \deleted{\unsure{disons|fixé|dis\textsuperscript{t}}} $i<2d$

d'où également

\[
\underline H^i_Y(\unsure{\mu_{nX}|\mu_{n_X}}^{\otimes d})=0\qquad\text{pour }i<2d
\]

C'est la situation du th. de pureté, vrai

(i) $X,Y$ lisses sur une base $S$

(ii) $Y,\unsure{\deleted{\to}X|\to X}$ réguliers, $X$ excellent, car. nulle

(iii) $d=1$, $X$ normal.

% archive sheet number: 25
\scanpage{27}
\mspage{- 216 -}\hfill n\textsuperscript{o} 391

(\S~8, n\textsuperscript{o}~3, formule (2)) ; autrement dit, la suite des exposants est :

\[
1,\ 5,\ 7,\ 9,\ 11,\ 13,\ 17.
\]

(X) On déduit de (IX) et du \S~8, n\textsuperscript{o}~3, cor. 1 de la prop. 3, que l'ordre du groupe de Weyl $W(R)$ est :

\[
2.6.8.10.12.14.18=2^{10}.3^4.5.7.
\]

(XI) Le même raisonnement qu'au n\textsuperscript{o}~8 montre que le seul automorphisme du schéma de Dynkin $E_7$ est l'identité.

12. \uline{Système de type} $E_6$.

(I) Gardons toujours les notations du n\textsuperscript{o}~10, et considérons l'intersection $R''$ de $R$ et du sous-espace $V''$ de codimension 2, orthogonal aux deux vecteurs $\varepsilon_8+\varepsilon_7$ et $\varepsilon_7-\varepsilon_6$ ; ce sont donc les vecteurs $x=\displaystyle\sum_{i=1}^{8}\xi_i\varepsilon_i$ de $R$ tels que $\xi_8=-\xi_7=-\xi_6$, autrement dit :

\[
\begin{aligned}
&\pm\varepsilon_i\pm\varepsilon_j,\quad(i\neq j,\ 1\leq i\leq5,\ 1\leq j\leq5)\\
&\pm\frac{1}{2}\bigl(\varepsilon_8-\varepsilon_7-\varepsilon_6+\sum_{i=1}^{5}(-1)^{\unsure{\nu|\partial|\mathcal D}(i)}\varepsilon_i\bigr)
\quad\text{où la somme }\sum_{i=1}^{5}\nu(i)\text{ est}
\end{aligned}
\]
\uline{paire}.

Il est clair que $R''$ engendre $V''$, donc est un système de racines réduit de rang 6, ayant $n=\binom{5}{2}.4+2^5=72$ éléments.

(II) Si l'on exprime qu'une racine $\displaystyle\sum_{i=1}^{8}n_i\alpha_i$ de $R$ est orthogonale à la fois à $\varepsilon_8+\varepsilon_7$ et à $\varepsilon_7-\varepsilon_6$, il vient, en vertu du n\textsuperscript{o}~10, (II), $n_8=n_7=0$. Le même raisonnement que dans le n\textsuperscript{o}~11 montre alors que les $\alpha_i$ pour $1\leq i\leq6$ forment un système de \uline{racines simples} $S''$ de $R''$. Le calcul des $(\alpha_i\mid\alpha_j)$ fait au n\textsuperscript{o}~10 montre que le schéma de Dynkin de $R''$ est de type $E_6$, et $R''$ est irréductible.

(III) En vertu du \S~8, n\textsuperscript{o}~3, formule (4), on a $h=n/\ell=12$.

% archive sheet number: 26
\scanpage{28}
Alors la suite \unsure{spectral|spectrale} implique

\[
\begin{aligned}
H^i_Y(X,\unsure{\mu_{nX}|\mu_{n_X}}^{\otimes d})&=0\qquad\text{si }i<2d\\
H^{2d}_Y(X,\mu_{nX}^{\otimes d})&\simeq\Gamma(Y,\underline H^{2d}_Y(\deleted{\unsure{\mathbb Z|Z|\mathbb Z/n\mathbb Z|\mathbb Z/n}}\mu_{nX}^{\otimes d}))
\end{aligned}
\]

d'où un \deleted{\unsure{lieu|hom|base|bon} canonique} élément canonique

\[
\gamma^{(n)}_{X/Y}\in H^{2d}_Y(X,\mu_{nX}^{\otimes d})
\]

[qu'on peut maintenant spécialiser dans $H^{2d}_{\deleted{\unsure{Y|Z}}}(X,\mu_{nX}^{\otimes d})$, ou si $X$ est propre sur un $S$, dans

\[
H^{2d}_{\Phi_{Y/X}}(X,\mu_{nX}^{\otimes d}).
\]

\margin{$\Phi\longrightarrow$ famille de \unsure{spts|cpts}}

\margin{N.B \inserted{\unsure{$d=1$|si|div.}} $\beta_{Y/X}\in H^1_Y(X,G_{mX})$ et $\gamma^{(n)}_{Y/X}=\partial_n(\beta_{Y/X})$, en particulier son image dans $H^2_Y(X,G_{mX})$ est nulle}

\uline{Exemple 1}. $X$ lisse sur $k$ \inserted{sép. clos.}, \deleted{$x$ pt} dim relative $d$, \deleted{il va résulter des le th. de} \inserted{$x$ pt de $X$ rat sur $k$, d'où}

\[
\gamma_{x/X}\in H^{2d}_!(X,\mu_{nX}^{\otimes d}).
\]

Il va résulter de la théorie qu'on va rappeler que \deleted{pour \unsure{tt}} si $X$ est géom. connexe, cet élément est \uline{une base} \uline{de} $H^{2d}_!(X,\mu_{nX}^{\otimes d})$ \uline{et ne dépend pas du choix de $x$}\unsure{;|.} \deleted{\unsure{ce|que} \illegible{1}} C'est qu'on avait vérifié directement dans le cas où $d=1$, en \deleted{\unsure{le|en|les|la} \unsure{compa|compar|corps}}\inserted{nous ramenant} au cas où $C$ était complète et où on dispose \deleted{\unsure{déjà}}

% archive sheet number: 27
\scanpage{29}
\mspage{- 214 -}\hfill n\textsuperscript{o} 391

son ordre est $h-1=17$. Elle est orthogonale à toutes les racines simples sauf à $\alpha_1$ et l'on a $(\tilde\alpha\mid\alpha_1)=1$.

(V) Ce qui a été vu au n\textsuperscript{o}~10 montre que les \uline{coracines} de $R'$ s'obtiennent encore en remplaçant $\varepsilon_i$ par $\varepsilon_i^*$ dans les expressions des racines de $R'$, et $R'{}^*$ est donc isomorphe à $R'$. Toutes les racines ayant même longueur, la \uline{forme bilinéaire canonique} pour $R$ est $(x\mid y)/2h=(x\mid y)/36$ et la longueur des racines est $1/3\sqrt{2}$, (\S~8, n\textsuperscript{o}~5, \uline{Remarque}). Le nombre des racines non orthogonales à une racine donnée est $4h-6=66$ et le nombre $\gamma$ est égal à $h^2=2^2.3^4$.

(VI) Le calcul des poids fondamentaux, formant la base duale de $(\alpha_i^*)$, $(1\leq i\leq7)$, dans $V'$, donne :

\[
\begin{aligned}
\omega_1&=\varepsilon_8-\varepsilon_7=2\alpha_1+2\alpha_2+3\alpha_3+4\alpha_4+3\alpha_5+2\alpha_6+\alpha_7=\tilde\alpha\\
\omega_2&=\frac{1}{2}(\varepsilon_1+\varepsilon_2+\varepsilon_3+\varepsilon_4+\varepsilon_5+\varepsilon_6-2\varepsilon_7+2\varepsilon_8)=\\
&=\frac{1}{2}(4\alpha_1+7\alpha_2+8\alpha_3+12\alpha_4+9\alpha_5+6\alpha_6+3\alpha_7)\\
\omega_3&=\frac{1}{2}(-\varepsilon_1+\varepsilon_2+\varepsilon_3+\varepsilon_4+\varepsilon_5+\varepsilon_6-3\varepsilon_7+3\varepsilon_8)=\\
&=3\alpha_1+4\alpha_2+6\alpha_3+8\alpha_4+6\alpha_5+4\alpha_6+2\alpha_7\\
\omega_4&=\varepsilon_3+\varepsilon_4+\varepsilon_5+\varepsilon_6+2(\varepsilon_8-\varepsilon_7)=\\
&=4\alpha_1+6\alpha_2+8\alpha_3+12\alpha_4+9\alpha_5+6\alpha_6+3\alpha_7\\
\omega_5&=\frac{1}{2}(2\varepsilon_4+2\varepsilon_5+2\varepsilon_6+3(\varepsilon_8-\varepsilon_7))=\\
&=\frac{1}{2}(6\alpha_1+9\alpha_2+12\alpha_3+18\alpha_4+15\alpha_5+10\alpha_6+5\alpha_7)\\
\omega_6&=\varepsilon_5+\varepsilon_6-\varepsilon_7+\varepsilon_8=2\alpha_1+3\alpha_2+4\alpha_3+6\alpha_4+5\alpha_5+4\alpha_6+2\alpha_7\\
\omega_7&=\varepsilon_6+\frac{1}{2}(\varepsilon_8-\varepsilon_7)=\frac{1}{2}(2\alpha_1+3\alpha_2+4\alpha_3+6\alpha_4+5\alpha_5+4\alpha_6+3\alpha_7).
\end{aligned}
\]

% archive sheet number: 28
\scanpage{30}
(par la th. de Kummer \unsure{globale|global}) d'un isom canonique

\[
t:H^2(\deleted{\unsure{C}}X,\unsure{\mu_{nX}|\mu_{n_X}})\xrightarrow{\sim}\mathbb Z/n\mathbb Z.
\]

On \unsure{avait|aurait} vérifié simplement

\[
t(\gamma_{x/X})=1
\]

\uline{Exemple} 2 \quad $X$ et $Y$ lisses sur $S$, $Y$ de codim $1$ dans $X$ en chaque pt. Alors on est \unsure{sur|sous} les conditions favorables, qui impliquent qu'on a aussi un hom

\diagram{\begin{tikzcd}[ampersand replacement=\&,column sep=large]
\Gamma(Y,\underline H^{2d}_Y(\mu_{nX}^{\otimes d}))\arrow[d,dash,"\wr"']\arrow[r,draw=none,"{\deleted{\longrightarrow}}" description] \& \deleted{H^0(S,R^{2d}_! f_*(\mu}\\
H^{2d}_Y(X,\mu_{nX}^{\otimes d})\arrow[r] \& \deleted{\unsure{H|\Gamma}}\Gamma(S,\mathbb R^{2d}_! f_*(\mu_{nX}^{\otimes d}))
\end{tikzcd}}

% archive sheet number: 29
\scanpage{31}
\mspage{- 213 -}\hfill n\textsuperscript{o} 391

11. \uline{Système de type} $E_7$\unsure{.|:}

(I) Gardons les notations du n\textsuperscript{o}~10, et considérons l'intersection $R'$ de $R$ et de l'hyperplan $V'$ de $V$ orthogonal au vecteur $\varepsilon_7+\varepsilon_8$ ; ce sont donc les vecteurs $x=\displaystyle\sum_{i=1}^{8}\xi_i\varepsilon_i$ de $R$ tels que $\xi_7+\xi_8=0$, autrement dit :

\[
\begin{aligned}
&\pm\varepsilon_i\pm\varepsilon_j,\quad(i\neq j,\ 1\leq i\leq6,\ 1\leq j\leq6),\ \pm(\varepsilon_7-\varepsilon_8)\text{ et}\\
&\pm\frac{1}{2}\bigl(\varepsilon_7-\varepsilon_8+\sum_{i=1}^{6}(-1)^{\unsure{\nu|\partial|\mathcal D}(i)}\varepsilon_i\bigr),\quad\text{où }\sum_{i=1}^{6}\nu(i)\text{ est \uline{impair}.}
\end{aligned}
\]

Il est clair que $R'$ engendre $V'$, donc est un système de racines réduit de rang 7, ayant $n=2+\binom{6}{2}.4+2^6=126$ éléments.

(II) Munissons $V'$ de l'ordre induit par l'ordre lexicographique sur $V$ ; les racines de $R'$ qui sont positives pour cet ordre sont, comme on l'a vu au n\textsuperscript{o}~10, de la forme $\displaystyle\sum_{i=1}^{8}n_i\alpha_i$, où les $n_i$ sont des entiers $\geq0$ ; si l'on exprime qu'une telle racine est orthogonale à $\varepsilon_7+\varepsilon_8$, il vient, (puisque la plus grande racine $\varepsilon_7+\varepsilon_8$ de $R$ est orthogonale à tous les $\alpha_i$ sauf $\alpha_8$), $n_8=0$. Toutes les racines positives de $R'$, (pour l'ordre considéré), sont donc combinaisons linéaires à coefficients positifs des $\alpha_i$ pour $1\leq i\leq7$, ce qui montre que ces racines, (qui sont orthogonales à $\varepsilon_7+\varepsilon_8$, donc appartiennent à $R'$), forment un système de \uline{racines simples} $S'$ de $R'$. Le calcul des $(\alpha_i\mid\alpha_j)$ fait au n\textsuperscript{o}~10 montre alors que le schéma de Dynkin de $R'$ est de type $E_7$, et $R'$ est irréductible.

(III) En vertu du \S~8, n\textsuperscript{o}~3, formule (4), on a $h=n/\ell=18$.

(IV) La plus grande racine pour l'ordre induit par celui de $R^8$ défini au n\textsuperscript{o}~10 est maintenant $\varepsilon_8-\varepsilon_7=2\alpha_1+2\alpha_2+3\alpha_3+4\alpha_4+3\alpha_5+2\alpha_6+\alpha_7$, donc c'est la plus grande racine $\tilde\alpha$, (\S~6, n\textsuperscript{o}~6, cor. 2 de la prop. 13), et on vérifie que

% archive sheet number: 30
\scanpage{32}
1\deleted{\unsure{5}}6) \uline{Classes fondamentales} \uline{globales}

On fixe un entier $n>0$, et on se \deleted{Soit $f:X\unsure{\mathrel{-}|\longrightarrow}Y$ morphisme} borne à des faisceaux de $n$-torsion (i.e. Modules sur $A=\mathbb Z/n\mathbb Z$), et des préschémas à \unsure{car.} résiduelles premières à $n$ i.e tels que $n.1_{\mathcal O_X}$ est inversible dans $\mathcal O_X$.

Si $f:X\mathrel{-}Y$ est un morphisme lisse, on dira qu'il est \inserted{purement} de dim relative $d$ si ses fibres sont \inserted{purement} de dim $d$. On \inserted{désigne} par $T^{(n)}_{X/Y}$ ou simplement $T_{X/Y}$ s'il n'y a pas de confusion \unsure{possible|possibles}, le faisceau $\unsure{\mu_{nX}|\mu_{n_X}}^{\otimes d}$. [Si $f$ est lisse de dim \inserted{relative \unsure{év\textsuperscript{t}.|év!} variable} quelconque, alors on définit de façon évidente $T_{X/Y}$ en le posant sur les divers morceaux].

Nous \unsure{regardons|regarderons} par la suite des morphismes \uline{lisses} \uline{compactifiables} \uline{purs} \uline{de dim} $d$ \qquad $f:X\longrightarrow Y$.

Pour un tel morphisme, on peut \unsure{définir|définit} un \uline{\deleted{\unsure{prem|prem.}}\inserted{hom.} canonique}

% archive sheet number: 31
\scanpage{33}
\mspage{- 188 -}\hfill n\textsuperscript{o} 391

On représentera un schéma de Dynkin $(X,f)$ par la figure qui représente son graphe sous-jacent $X$, (\S~3, n\textsuperscript{o}~1), chaque sommet étant surmonté de la valeur de $f$ en ce sommet. Lorsque $f$ est constante et égale à 1, on omet d'indiquer ses valeurs.

Il est clair que les schémas de Dynkin de $R$ pour deux chambres différentes sont isomorphes puisque le groupe de Weyl $W(R)$ permute les chambres.

En outre la connaissance du schéma de Dynkin d'un système de racines détermine ce dernier à isomorphisme près. De façon précise :

PROPOSITION 2.- \uline{Soient} $V$, $V'$ \uline{deux espaces vectoriels euclidiens}, $R$, (resp. $R'$), \uline{un système réduit de racines dans} $V$, (resp. $V'$), $(v_i)_{1\leq i\leq\ell}$, (resp. $(v'_j)_{1\leq j\leq\ell'}$), \uline{les racines simples de} $R$, (resp. $R'$), \uline{relatives à une chambre}. \uline{S'il existe un isomorphisme} $v_i\longrightarrow v'_{\varphi(i)}$ \uline{du schéma de Dynkin de} $R$ \uline{sur celui de} $R'$, \uline{il existe un isomorphisme de} $V$ \uline{sur} $V'$ \uline{transformant} $v_i$ \uline{en} $v'_{\varphi(i)}$ \uline{pour tout} $i$, \uline{et transformant} $R$ \uline{en} $R'$.

L'hypothèse entraîne $\ell=\ell'$ et, (après multiplication de la norme dans $V'$ par un facteur constant), $(v_i\mid v_k)=(v'_{\varphi(i)}\mid v'_{\varphi(k)})$ quels que soient $i$ et $k$. L'application $\psi:\xi_1v_1+..+\xi_\ell v_\ell\longrightarrow\xi_1v'_{\varphi(1)}+..+\xi_\ell v'_{\varphi(\ell)}$ est donc un isomorphisme de $V$ sur $V'$, qui transforme le groupe de Weyl $W(R)$ en le groupe de Weyl $W(R')$. Or $R$, (resp. $R'$), est l'ensemble des transformés des $v_i$, (resp. $v'_j$), par $W(R)$, (resp. $W(R')$), (\S~6, n\textsuperscript{o}~5, prop. 10), et $\psi$ transforme donc $R$ en $R'$.

COROLLAIRE.- \uline{Si} $A(R)$ \uline{est le sous-groupe de} $\mathrm{GL}(V)$ \uline{laissant stable un système de racines} $R$, \uline{le groupe quotient} $A(R)/W(R)$ \uline{est isomorphe au groupe des automorphismes du schéma de Dynkin de} $R$.

% archive sheet number: 32
